% [claude] 2026-11-16 draft. Protocol follows reports/2026-11-16-64-retrieval-baseline-matrix-plan.md.
\section{Problem Setup and Evaluation Protocol}
\label{sec:setup}

\paragraph{Task.}
A gallery $\mathcal{G}$ contains images, and a query $q$ is a sentence that states a set of conditions $C(q)=\{c_1,\dots,c_k\}$.
An image $x$ is relevant, $\rel(q,x)=1$, only if it satisfies every condition in $C(q)$.
Images whose status cannot be decided from the annotations form an ignore set $\mathcal{I}(q)$ and are removed from the ranking of $q$ (\cref{sec:benchmark}).
The positives are $\mathcal{P}(q)=\{x\in\mathcal{G}:\rel(q,x)=1\}$.

\paragraph{Two-stage retrieval.}
A first-stage model with score $s_f(q,x)$ ranks the whole gallery.
A reranker with score $s_r(q,x)$ re-scores the top-$K$ images of that ranking.
The final ranking places the reranked top-$K$ first, followed by the remaining images in their first-stage order.
With $K{=}0$ this reduces to embedding-only retrieval.

\paragraph{Metrics.}
For each query we compute average precision over the full ranking of $\mathcal{G}\setminus\mathcal{I}(q)$, with $|\mathcal{P}(q)|$ as the denominator, so positives outside the top-$K$ lower the score instead of being dropped; mAP averages over queries.
We also report Recall@$K$, the fraction of positives that the first stage places in its top $K$.
To separate the two failure modes of a two-stage system we use a fixed-candidate oracle: it moves every positive already in the top-$K$ to the front and keeps everything else unchanged.
The gap between this oracle and the system measures in-pool ordering errors that a better reranker could fix; the gap between a perfect ranking and the oracle measures positives that the first stage leaves outside the pool or ranks poorly in the tail.

\paragraph{Candidate depth.}
Our primary setting is $K{=}100$, matching the reranking evaluations of Qwen3-VL-Reranker and the CORE development set~\cite{qwen3vlemb2026,core2026}.
We vary $K\in\{50,200\}$ to study how depth trades recall for reranking cost.
The same $K$ does not imply the same compute: we report measured wall time and peak memory for every first stage and reranker.

\paragraph{Development and test discipline.}
Every benchmark has a development portion and a test portion.
We select methods, prompts, and hyperparameters only on development data.
Baselines are run without tuning, following each original paper, and we check each one against the numbers its authors report before including it (\cref{sec:experiments}).
Our own methods are evaluated on test data once, after all choices are fixed.
Where a benchmark releases only a test set, we use it only in this final evaluation.

\paragraph{Uncertainty.}
In our benchmark the positive sets of different queries are disjoint, so each query is an independent unit, and we report 95\% bootstrap intervals over queries.
For other benchmarks we follow their unit of analysis and resample groups of queries that share images.
